\VerdierPhysicalPageAnchor{0081}%
\VerdierSemanticAnchor{ch1-2-5-11}
\noindent\textbf{2.5.11.}\quad
\VerdierSemanticAnchor{ch1-2-5-11-homotopism}
Soient $X$ et $Y$ deux complexes de $C$. Un morphisme de
$\operatorname{comp}(C)$, $f\colon X\longrightarrow Y$, est appelé un
\emph{homotopisme} si son image $\dot f$ dans $K(C)$ est un isomorphisme.
\VerdierSemanticAnchor{ch1-2-5-11-homotopic-complexes}
Les complexes $X$ et $Y$ sont dits \emph{homotopes} s'ils sont isomorphes
dans $K(C)$.

\VerdierSemanticAnchor{ch1-2-5-12}
\noindent\textbf{2.5.12.}\quad
\VerdierSemanticAnchor{ch1-2-5-12-setup}
Soient $X$ et $Y$ deux complexes de $C$,
$f,g\colon X\rightrightarrows Y$ deux morphismes de
$\operatorname{comp}(C)$,
$a,b\colon f\rightsquigarrow g$ deux homotopies reliant $f$ à $g$.
\VerdierSemanticAnchor{ch1-2-5-12-second-order-homotopy}
On appelle \emph{homotopie du deuxième ordre} reliant $a$ à $b$, notée
$\alpha\colon a\rightsquigarrow b$, une famille de morphismes
\[
  \alpha=(\alpha^i\colon X^{i+2}\longrightarrow Y^i,\ i\in\mathbb Z)
\]
telle que :

\VerdierSemanticAnchor{ch1-2-5-12-eq1}
\begin{equation}
  b^i-a^i
  =\alpha^i d_X^{i+1}-d_Y^{i-1}\alpha^{i-1}.
  \tag{2.5.12.1}
\end{equation}

\VerdierSemanticAnchor{ch1-2-5-13}
\noindent\textbf{2.5.13.}\quad
Remarquons pour terminer que l'isomorphisme (2.4.1) :

\VerdierSemanticAnchor{ch1-2-5-13-complex-opposite-isomorphism}
\begin{equation*}
  \iota_1\colon
  \operatorname{comp}(C)^\circ
  \longrightarrow
  \operatorname{comp}(C^\circ)
\end{equation*}
définit, en vertu de (2.5.7), un unique isomorphisme :

\VerdierSemanticAnchor{ch1-2-5-13-homotopy-opposite-isomorphism}
\begin{equation*}
  \iota_1\colon
  K(C)^\circ
  \longrightarrow
  K(C^\circ)
\end{equation*}
\VerdierSemanticAnchor{ch1-2-5-13-z-category-isomorphism}
qui est un isomorphisme de $\mathbb Z$-catégories.
\VerdierSemanticAnchor{ch1-2-5-13-identification}
Nous identifierons dorénavant le plus souvent les catégories $K(C)^\circ$ et
$K(C^\circ)$ par cet isomorphisme.
